% ============================================================
% PHẦN 2: CÁC CÔNG TRÌNH NGHIÊN CỨU LIÊN QUAN
% ============================================================
\section{Các công trình nghiên cứu liên quan}

% ------------------------------------------------------------
\subsection{Phát hiện bất thường trên mạng cấp nước}

Phát hiện bất thường trên mạng cấp nước (Water Distribution Network --- WDN) là
một bài toán đã thu hút nhiều sự quan tâm trong cộng đồng nghiên cứu, đặc biệt kể
từ khi các hệ thống điều khiển và thu thập dữ liệu (SCADA) được tích hợp rộng rãi
vào cơ sở hạ tầng đô thị.

Nghiên cứu của Taormina và cộng sự \cite{taormina2018batadal} đã tổ chức cuộc thi
BATADAL (Battle of the Attack Detection ALgorithms) --- một benchmarking challenge
nhằm đánh giá khả năng phát hiện tấn công mạng vật lý (cyber-physical attack) trên
hệ thống WDN. Bộ dữ liệu BATADAL được xây dựng trên mô hình mạng nước C-TOWN với
43 cảm biến đo mức nước bể, lưu lượng bơm, áp suất và trạng thái van. Kết quả
cuộc thi cho thấy sự đa dạng trong cách tiếp cận: từ các phương pháp thống kê
truyền thống (statistical process control), đến các mô hình học máy giám sát
(supervised learning) như Random Forest và Support Vector Machine, lẫn phương pháp
dựa trên mô hình vật lý (model-based detection).

Eliades và Polycarpou \cite{eliades2012leakage} đề xuất phương pháp phát hiện rò
rỉ dựa trên mô hình số học của mạng nước, sử dụng bộ quan sát thích nghi
(adaptive observer) để ước lượng trạng thái hệ thống và phát hiện sai lệch so với
giá trị kỳ vọng. Phương pháp này có ưu điểm về khả năng giải thích (interpretability)
nhưng đòi hỏi một mô hình thủy lực chính xác, điều không dễ có được trong thực tế.

Nghiên cứu của Kravchuk và Polyanska \cite{kravchuk2019anomaly} so sánh nhiều thuật
toán học máy cho bài toán phát hiện bất thường trên dữ liệu SCADA của hệ thống
nước, bao gồm Isolation Forest, Local Outlier Factor và Autoencoder. Kết quả chỉ ra
rằng các phương pháp không giám sát (unsupervised) phù hợp hơn trong môi trường
thực tế khi nhãn tấn công rất hiếm, trong khi các phương pháp giám sát như Random
Forest đạt độ chính xác cao hơn khi dữ liệu có nhãn đầy đủ.

Soltan, Bassily và Matta \cite{soltan2018react} phát triển hệ thống REACT
(Real-Time Attack-Resilient Control) có khả năng phát hiện và giảm thiểu tấn công
vào hệ thống WDN theo thời gian thực. Đây là một trong những công trình tiên phong
kết hợp phát hiện bất thường với cơ chế phản ứng tự động, mở ra hướng nghiên cứu
về resilient control cho cơ sở hạ tầng quan trọng.

Trong khuôn khổ đề tài này, phương pháp rule-based anomaly detection được xây dựng
dựa trên thông số vận hành của mạng C-TOWN (ngưỡng mức nước bể, trạng thái bơm
bất thường) và được tích hợp song song với mô hình học máy Random Forest để so
sánh hiệu quả của hai cách tiếp cận trên cùng một luồng dữ liệu streaming.

% ------------------------------------------------------------
\subsection{Kiến trúc Big Data cho dữ liệu IoT}

Việc xử lý và phân tích dữ liệu từ các thiết bị IoT quy mô lớn đặt ra những yêu
cầu đặc thù về kiến trúc hệ thống, đặc biệt là khả năng xử lý luồng dữ liệu liên
tục với độ trễ thấp.

Marz và Warren \cite{marz2015big} đề xuất kiến trúc Lambda --- một trong những mô
hình kiến trúc Big Data có ảnh hưởng nhất --- với ba lớp độc lập: Batch Layer (xử
lý toàn bộ dữ liệu lịch sử theo lô), Speed Layer (xử lý dữ liệu mới nhất theo
thời gian thực) và Serving Layer (hợp nhất kết quả từ hai lớp trên để phục vụ
truy vấn). Ưu điểm của kiến trúc Lambda là tính chịu lỗi cao và đảm bảo tính nhất
quán của kết quả. Tuy nhiên, hạn chế lớn nhất là sự trùng lặp trong logic xử lý
giữa Batch Layer và Speed Layer, dẫn đến chi phí vận hành và bảo trì cao.

Kreps \cite{kreps2014kappa} đề xuất kiến trúc Kappa như một giải pháp thay thế
tinh gọn hơn: loại bỏ hoàn toàn Batch Layer và chỉ duy trì một Speed Layer duy
nhất dựa trên hệ thống xử lý luồng. Trong kiến trúc này, Apache Kafka đóng vai trò
trung tâm như một distributed log bền vững, cho phép ``replay'' toàn bộ lịch sử dữ
liệu khi cần tái xử lý (reprocessing). Kiến trúc Kappa được ứng dụng rộng rãi cho
các hệ thống IoT nơi dữ liệu liên tục phát sinh và yêu cầu xử lý tức thì.

Kafka \cite{kreps2011kafka} là hệ thống message broker phân tán, được thiết kế để
xử lý luồng sự kiện (event stream) với thông lượng cao và độ trễ thấp. Kafka sử
dụng mô hình publish-subscribe với topic và partition, đảm bảo thứ tự thông điệp
trong từng partition và cho phép nhiều consumer đọc dữ liệu độc lập. Khả năng lưu
trữ dữ liệu lâu dài (retention) và cơ chế offset tracking làm cho Kafka trở thành
lựa chọn phổ biến trong các pipeline IoT.

Zaharia và cộng sự \cite{zaharia2016spark} giới thiệu Apache Spark như một unified
analytics engine cho Big Data, với khả năng xử lý cả batch và streaming trên cùng
một nền tảng. Spark Structured Streaming, được giới thiệu trong Spark 2.0, cung
cấp mô hình lập trình thống nhất dựa trên DataFrame/Dataset, cho phép viết logic
xử lý streaming giống hệt batch processing. Micro-batch execution với trigger
interval linh hoạt cho phép điều chỉnh cân bằng giữa throughput và latency.

Kamburugamuve và Fox \cite{kamburugamuve2016survey} thực hiện khảo sát toàn diện
các hệ thống xử lý luồng dữ liệu IoT, so sánh Apache Storm, Spark Streaming,
Apache Flink và các nền tảng thương mại. Kết quả cho thấy Spark Structured Streaming
nổi bật về tính dễ sử dụng và tích hợp tốt với hệ sinh thái Hadoop/HDFS, trong khi
Flink có ưu thế về true streaming (xử lý từng sự kiện) và event-time semantics.

Đề tài này áp dụng kiến trúc Kappa với Kafka làm distributed log và Spark Structured
Streaming làm Speed Layer, phù hợp với đặc thù của bài toán: dữ liệu cảm biến WDN
phát sinh đều đặn theo chu kỳ 1 giờ, không yêu cầu xử lý lịch sử phức tạp, và cần
phản hồi cảnh báo trong thời gian thực.

% ------------------------------------------------------------
\subsection{Hệ thống giám sát thời gian thực}

Giám sát hệ thống phân tán theo thời gian thực là một bài toán kỹ thuật quan trọng
trong nhiều lĩnh vực: từ giám sát cơ sở hạ tầng CNTT, mạng công nghiệp
(Industrial IoT), đến theo dõi sức khỏe cơ sở hạ tầng đô thị.

Bernstein \cite{bernstein2014containers} nhận định rằng container hóa ứng dụng
(containerization) bằng Docker là xu hướng tất yếu trong triển khai hệ thống phân
tán hiện đại, giúp đảm bảo tính tái lập (reproducibility) giữa môi trường phát
triển và sản xuất. Docker Compose cho phép định nghĩa và khởi động toàn bộ stack
đa dịch vụ bằng một file cấu hình duy nhất, đặc biệt phù hợp cho các pipeline
Big Data với nhiều thành phần phụ thuộc lẫn nhau.

InfluxDB \cite{influxdata2024influxdb} là cơ sở dữ liệu chuỗi thời gian
(time-series database) được tối ưu cho việc ghi và truy vấn dữ liệu theo timestamp
với tần suất cao. Mô hình dữ liệu của InfluxDB gồm measurement, tag và field cho
phép truy vấn linh hoạt theo nhiều chiều. Ngôn ngữ truy vấn Flux, được giới thiệu
từ InfluxDB 2.x, cung cấp khả năng biến đổi và phân tích dữ liệu chuỗi thời gian
mạnh mẽ với cú pháp pipeline tương tự dplyr trong R. InfluxDB được sử dụng rộng
rãi trong các hệ thống giám sát IoT, monitoring cơ sở hạ tầng và observability.

Grafana \cite{grafana2024grafana} là nền tảng trực quan hóa và quan sát
(observability platform) mã nguồn mở, hỗ trợ hàng chục nguồn dữ liệu khác nhau bao
gồm InfluxDB, Prometheus, Elasticsearch và nhiều RDBMS. Tính năng alerting tích hợp
cho phép cấu hình cảnh báo theo ngưỡng (threshold-based) hoặc theo biểu thức phức
tạp, gửi thông báo qua email, Slack hoặc PagerDuty. Grafana Provisioning --- cơ chế
cấu hình tự động qua file YAML --- cho phép phiên bản hóa (version control) toàn
bộ cấu hình dashboard và datasource, rất phù hợp với các hệ thống cần triển khai
tự động (Infrastructure as Code).

Nakagawa và cộng sự \cite{nakagawa2020realtime} xây dựng hệ thống giám sát
cơ sở hạ tầng nước đô thị sử dụng stack Kafka --- InfluxDB --- Grafana, xử lý dữ
liệu từ hơn 200 cảm biến IoT phân bố trên toàn mạng. Nghiên cứu cho thấy kiến
trúc này đáp ứng yêu cầu end-to-end latency dưới 500\,ms cho pipeline từ cảm biến
đến dashboard, đồng thời dễ dàng mở rộng khi số lượng cảm biến tăng lên.

Chandola, Banerjee và Kumar \cite{chandola2009anomaly} tổng hợp toàn diện các kỹ
thuật phát hiện bất thường, phân loại thành ba nhóm: point anomaly (một điểm dữ
liệu lệch khỏi phân phối bình thường), contextual anomaly (điểm dữ liệu bất thường
trong ngữ cảnh cụ thể) và collective anomaly (tập hợp điểm dữ liệu tạo thành mẫu
bất thường). Công trình này là tài liệu tham khảo nền tảng cho hầu hết các nghiên
cứu về anomaly detection trong chuỗi thời gian, bao gồm dữ liệu cảm biến WDN.

Tổng kết lại, đề tài kế thừa các kết quả nghiên cứu hiện có theo ba hướng: (1)
sử dụng bộ dữ liệu chuẩn BATADAL \cite{taormina2018batadal} làm cơ sở đánh giá
khách quan; (2) áp dụng kiến trúc Kappa \cite{kreps2014kappa} với Kafka
\cite{kreps2011kafka} và Spark Structured Streaming \cite{zaharia2016spark} làm
nền tảng xử lý luồng dữ liệu; và (3) tích hợp InfluxDB cùng Grafana làm lớp lưu
trữ và trực quan hóa thời gian thực. Điểm khác biệt của đề tài so với các công
trình liên quan là sự kết hợp đồng thời của cả ba yếu tố này trong một pipeline
hoàn chỉnh, được containerize, kiểm thử toàn diện và triển khai thực tế trên nền
tảng đám mây AWS EC2.
